\documentclass[aspectratio=169,spanish,xcolor=table]{beamer}
\usepackage{iftex}
\ifPDFTeX
  \usepackage[utf8]{inputenc}
  \usepackage[T1]{fontenc}
\fi
\usepackage[spanish,es-tabla,es-noshorthands]{babel}
\usepackage{helvet}
\renewcommand{\familydefault}{\sfdefault}
\usepackage{tikz,pgfplots,array,booktabs,ragged2e}
\pgfplotsset{compat=1.18}
\usetikzlibrary{arrows.meta,positioning,calc,shapes.geometric,fit}

\definecolor{iiiepeNavy}{HTML}{163B5C}
\definecolor{iiiepeBlue}{HTML}{1F6F9F}
\definecolor{iiiepeAqua}{HTML}{20B7A8}
\definecolor{iiiepeOrange}{HTML}{F28C28}
\definecolor{iiiepePaper}{HTML}{F6F8F7}
\definecolor{iiiepeInk}{HTML}{1E2A32}
\definecolor{iiiepePale}{HTML}{EEF5F8}
\definecolor{iiiepeMint}{HTML}{E7F6F4}
\definecolor{iiiepeWarm}{HTML}{FFF3E6}

\usetheme{default}
\usecolortheme{default}
\setbeamertemplate{navigation symbols}{}
\setbeamercolor{normal text}{fg=iiiepeInk,bg=white}
\setbeamercolor{structure}{fg=iiiepeNavy}
\setbeamercolor{frametitle}{fg=white,bg=iiiepeNavy}
\setbeamercolor{block title}{fg=white,bg=iiiepeBlue}
\setbeamercolor{block body}{fg=iiiepeInk,bg=iiiepePale}
\setbeamercolor{item}{fg=iiiepeOrange}

\setbeamertemplate{frametitle}{%
  \nointerlineskip%
  \begin{beamercolorbox}[wd=\paperwidth,ht=0.95cm,dp=0.22cm,leftskip=0.45cm,rightskip=0.45cm]{frametitle}%
    \begin{minipage}[c]{13.2cm}\usebeamerfont{frametitle}\large\bfseries\insertframetitle\end{minipage}%
    \hfill\begin{minipage}[c]{1.6cm}\raggedleft\includegraphics[height=0.56cm]{img/iiiepe/logo-iiiepe-blanco.png}\end{minipage}%
  \end{beamercolorbox}%
}

\setbeamertemplate{footline}{%
  \leavevmode%
  \hbox{%
    \begin{beamercolorbox}[wd=.78\paperwidth,ht=0.5cm,dp=0.18cm,leftskip=0.45cm]{frametitle}%
      \scriptsize\color{white}M. de la Cruz Muñoz · MAEM · Matrícula 00874 · SD26 · Actividad 4 · Inv.-Acción%
    \end{beamercolorbox}%
    \begin{beamercolorbox}[wd=.22\paperwidth,ht=0.5cm,dp=0.18cm,rightskip=0.45cm]{frametitle}%
      \hfill\scriptsize\color{white}\insertframenumber{} / \inserttotalframenumber%
    \end{beamercolorbox}%
  }%
  \vskip0pt%
}

\newcommand{\Pill}[2]{%
  \tikz[baseline=(x.base)]{\node[rounded corners=3pt,fill=#1!14,draw=#1!75!black,thick,inner xsep=6pt,inner ysep=4pt,align=center] (x) {\scriptsize\bfseries #2};}%
}

\title{Investigación-Acción y Transformación Docente}
\subtitle{Actividad 4 · La Investigación-Acción como Estrategia para Transformar la Práctica}
\author{Mtro. Martín Jonathan de la Cruz Muñoz · Matrícula 00874}
\institute{IIIEPE · Maestría en Aprendizaje y Enseñanza de las Matemáticas · Grupo A}
\date{19 de noviembre de 2026}

\begin{document}

\begin{frame}[plain]
  \begin{tikzpicture}[remember picture,overlay]
    \fill[iiiepeNavy] (current page.north west) rectangle ([xshift=10.3cm]current page.south west);
    \fill[iiiepeBlue] ([xshift=10.3cm]current page.north west) rectangle (current page.south east);
    \fill[iiiepeOrange] ([xshift=0.8cm,yshift=-6.45cm]current page.north west) rectangle ([xshift=9.3cm,yshift=-6.53cm]current page.north west);
    \node[anchor=north west,text width=8.7cm,align=left] at ([xshift=0.8cm,yshift=-1.15cm]current page.north west) {%
      {\color{iiiepeAqua!45!white}\small\bfseries MAESTRÍA EN APRENDIZAJE Y ENSEÑANZA DE LAS MATEMÁTICAS · SD26\par}%
      \vspace{0.35cm}
      {\color{white}\Large\bfseries \inserttitle\par}%
      \vspace{0.25cm}
      {\color{white!90}\normalsize \insertsubtitle\par}%
    };
    \node[anchor=south west,text width=8.7cm,align=left] at ([xshift=0.8cm,yshift=0.75cm]current page.south west) {%
      {\color{white}\normalsize\bfseries Mtro. Martín Jonathan de la Cruz Muñoz\par}%
      {\color{white!85}\scriptsize Matrícula 00874 · martin.delacruz@campus.iiiepe.edu.mx\par}%
      {\color{iiiepeAqua!35!white}\scriptsize Monterrey, Nuevo León · 19 de noviembre de 2026\par}%
    };
    \node[anchor=center,text width=5.0cm,align=center] at ([xshift=13.15cm,yshift=-4.5cm]current page.north west) {%
      \includegraphics[width=3.6cm]{img/iiiepe/logo-iiiepe-blanco.png}\par\vspace{0.45cm}%
      {\color{white}\large\bfseries IIIEPE\par}%
      \vspace{0.1cm}%
      {\color{white!85}\scriptsize Fundamentos para la enseñanza\\y el aprendizaje I\par}%
    };
  \end{tikzpicture}
\end{frame}
\begin{frame}{La Investigación-Acción como Saber Pedagógico}
  \begin{block}{Tesis metodológica (Restrepo Gómez, 2004)}
    La investigación-acción educativa convierte el aula en un espacio de indagación reflexiva donde el docente deconstruye rutinas, planifica mejoras situadas, recoge evidencias y reconstruye su práctica con rigor.
  \end{block}
  \vspace{0.3cm}
  \begin{columns}[T]
    \column{0.48\textwidth}
      \textbf{\color{iiiepeNavy}¿Qué la distingue?}
      \begin{itemize}\scriptsize
        \item No separa investigador de practicante.
        \item Persigue transformar la enseñanza real, no solo describirla.
        \item Integra ciclos abiertos de acción y crítica.
      \end{itemize}
    \column{0.48\textwidth}
      \textbf{\color{iiiepeBlue}¿Por qué es clave en matemáticas?}
      \begin{itemize}\scriptsize
        \item Evita atribuir errores únicamente al estudiante.
        \item Permite rastrear obstáculos conceptuales específicos.
        \item Sostiene innovaciones verificables con evidencia.
      \end{itemize}
  \end{columns}
\end{frame}

\begin{frame}{Cinco Preguntas Fundamentales (1--3)}
  \begin{itemize}\small
    \item \textbf{\color{iiiepeNavy}1. Concepto y finalidad:} Indagación autorreflexiva orientada a transformar la práctica docente y mejorar los aprendizajes en su contexto real.
    \vspace{0.22cm}
    \item \textbf{\color{iiiepeBlue}2. Docente como investigador:} Quien enseña observa interacciones irrepetibles y puede contrastar sus supuestos con evidencias concretas de aula.
    \vspace{0.22cm}
    \item \textbf{\color{iiiepeAqua!75!black}3. Relación entre fases:} Deconstrucción/identificación $\rightarrow$ planificación $\rightarrow$ acción $\rightarrow$ observación $\rightarrow$ reflexión $\rightarrow$ nuevo ciclo mejorado.
  \end{itemize}
\end{frame}

\begin{frame}{Cinco Preguntas Fundamentales (4--5)}
  \begin{columns}[T]
    \column{0.48\textwidth}
      \begin{block}{4. Evidencias, observación y reflexión}
        \scriptsize Sin cuadernos, diarios de clase, listas de cotejo ni registros de diálogo, la mejora queda en impresión subjetiva. La evidencia muestra exactamente dónde se fractura la comprensión matemática.
      \end{block}
    \column{0.48\textwidth}
      \begin{block}{5. Alternativa de transformación}
        \scriptsize Rompe la repetición mecánica de programas año con año al instalar una cultura de ensayo pedagógico responsable, diálogo colegiado y mejora continua.
      \end{block}
  \end{columns}
\end{frame}

\begin{frame}{Espiral Helicoidal de la Investigación-Acción}
  \centering
  \begin{tikzpicture}[font=\scriptsize, >=Latex, scale=0.88, transform shape]
    \node[draw=iiiepeNavy,fill=iiiepePale,thick,circle,minimum size=2.1cm,text width=1.8cm,align=center,font=\scriptsize\bfseries,text=iiiepeNavy] (c) at (0,0) {Saber\\Pedagógico};
    \node[draw=iiiepeNavy,fill=white,thick,rounded corners=3pt,text width=3.4cm,align=center] (a) at (0,2.5) {\textbf{1. Identificación}\\Deconstruir con evidencias};
    \node[draw=iiiepeBlue,fill=white,thick,rounded corners=3pt,text width=3.4cm,align=center] (b) at (4.3,0.9) {\textbf{2. Planificación}\\Diseñar intervención};
    \node[draw=iiiepeAqua!80!black,fill=iiiepeMint,thick,rounded corners=3pt,text width=3.4cm,align=center] (d) at (2.8,-2.1) {\textbf{3. Acción}\\Aplicar en el grupo};
    \node[draw=iiiepeOrange!90!black,fill=iiiepeWarm,thick,rounded corners=3pt,text width=3.4cm,align=center] (e) at (-2.8,-2.1) {\textbf{4. Observación}\\Registrar producciones};
    \node[draw=iiiepeNavy,fill=iiiepePale,thick,rounded corners=3pt,text width=3.4cm,align=center] (f) at (-4.3,0.9) {\textbf{5. Reflexión}\\Decidir nuevo ciclo};
    \draw[->,very thick,iiiepeNavy] (a) -- (b);
    \draw[->,very thick,iiiepeBlue] (b) -- (d);
    \draw[->,very thick,iiiepeAqua!80!black] (d) -- (e);
    \draw[->,very thick,iiiepeOrange!90!black] (e) -- (f);
    \draw[->,very thick,iiiepeNavy] (f) -- (a);
  \end{tikzpicture}
\end{frame}

\begin{frame}{Problemática Delimitada en Matemáticas}
  \begin{block}{Foco observado en bachillerato técnico e ingeniería}
    Manipulación algebraica mecánica de ecuaciones y modelos lineales sin interpretar el significado físico de las variables, unidades ni límites de validez.
  \end{block}
  \vspace{0.25cm}
  \begin{columns}[T]
    \column{0.48\textwidth}
      \textbf{\color{iiiepeNavy}Evidencias iniciales:}
      \begin{itemize}\scriptsize
        \item Alumnado calcula pendientes correctamente de modo abstracto pero no interpreta unidades ($\text{m/s}$ o $\text{\$/h}$).
        \item Confusión reiterada entre modelo proporcional $y=kx$ y afín $y=mx+b$.
      \end{itemize}
    \column{0.48\textwidth}
      \textbf{\color{iiiepeOrange!85!black}Autocrítica docente:}
      \begin{itemize}\scriptsize
        \item Evitar culpar únicamente al rezago estudiantil.
        \item Reconocer que el exceso de ejercicios descontextualizados indujo dicho hábito.
      \end{itemize}
  \end{columns}
\end{frame}

\begin{frame}{Abordaje por Fases: Identificación, Planificación y Acción}
  \scriptsize\renewcommand{\arraystretch}{1.38}
  \begin{tabular}{p{2.5cm}p{5.5cm}p{5.0cm}}
    \toprule
    \rowcolor{iiiepePale} \textbf{Fase} & \textbf{Estrategia diseñada} & \textbf{Evidencia asociada} \\
    \midrule
    1. Identificación & Revisión diagnóstica de producciones sobre función lineal y significado físico. & Muestreo de cuadernos y errores de interpretación dimensional. \\
    2. Planificación & Secuencia de tres sesiones basada en tablas, gráficas y justificación escrita de parámetros. & Guía de retos contextualizados y matriz de valoración formativa. \\
    3. Acción & Trabajo en binas con predicción previa al cálculo y discusión plenaria de supuestos. & Hojas de trabajo por equipo y registro cronológico de aula. \\
    \bottomrule
  \end{tabular}
\end{frame}

\begin{frame}{Abordaje por Fases: Observación y Reflexión}
  \begin{columns}[T]
    \column{0.48\textwidth}
      \begin{block}{4. Observación sistemática}
        \begin{itemize}\scriptsize
          \item Diario docente enfocado en preguntas que desbloquean la argumentación física.
          \item Lista de cotejo sobre coherencia entre tabla, gráfica, ecuación y unidades.
          \item Contraste de explicaciones iniciales frente a versiones revisadas.
        \end{itemize}
      \end{block}
    \column{0.48\textwidth}
      \begin{block}{5. Reflexión evaluativa}
        \begin{itemize}\scriptsize
          \item Evaluar en qué medida disminuyeron respuestas carentes de sentido físico.
          \item Ajustar el grado de andamiaje escrito para el segundo ciclo.
          \item Socializar hallazgos con academia de matemáticas del plantel.
        \end{itemize}
      \end{block}
  \end{columns}
\end{frame}

\begin{frame}{Por Qué Sostiene la Innovación Educativa}
  \centering
  \Pill{iiiepeNavy}{Parte de problemas reales} \hspace{0.2cm}
  \Pill{iiiepeBlue}{Articula teoría y evidencia} \hspace{0.2cm}
  \Pill{iiiepeAqua!75!black}{Empodera al colectivo docente}
  \vspace{0.45cm}
  \begin{block}{Argumento integrador}
    \small Innovar sin investigar la propia aula degenera en adopción superficial de modas. La investigación-acción garantiza que cada cambio didáctico tenga sentido contextual, control metodológico y compromiso ético con el aprendizaje de todos.
  \end{block}
\end{frame}

\begin{frame}{Conclusiones y Proyección Metodológica}
  \begin{itemize}\normalsize
    \item La investigación-acción convierte al docente en autor crítico de saber pedagógico \textbf{(Restrepo Gómez, 2004)}.
    \item Sus cinco fases mantienen una relación cíclica abierta orientada por evidencias observables.
    \item Aplicar este ciclo al problema de la interpretación de modelos lineales establece la base metodológica directa para promover el aprendizaje profundo en matemáticas.
  \end{itemize}
\end{frame}

\begin{frame}{Referencias Principales (APA 7)}
  \scriptsize\setlength{\parskip}{5pt}
  \begin{itemize}
    \item Boström, E., y Palm, T. (2023). Formative assessment practice in mathematics. \textit{Frontiers in Education}, 8, 1101192.
    \item Collin, J., y Quigley, A. (2021). \textit{Teacher feedback to improve pupil learning}. EEF.
    \item Henderson, P., et al. (2017). \textit{Improving mathematics in Key Stages 2 and 3}. EEF.
    \item Pimienta Prieto, J. H. (2012). \textit{Metodología constructivista: Guía para la planeación docente} (3.ª ed.). Pearson.
    \item Quiroga Aguilar, L., y Lara Órdenes, E. (2022). El aprendizaje profundo como herramienta para cambio en la visión de aprendizaje de una cultura escolar. \textit{REA}, 12(1).
    \item Restrepo Gómez, B. (2004). La investigación-acción educativa y la construcción de saber pedagógico. \textit{Educación y Educadores}, (7), 45--55.
    \item Star, J. R., et al. (2015). \textit{Teaching strategies for improving algebra knowledge}. IES.
  \end{itemize}
\end{frame}

\end{document}
